\section{Conclusions}\label{sec:concl}

We studied colour-neutral electroweak quirks with $\Lam\sim0.5$--$10\keV$ in a CMS Phase-2-like tracker, at truth level for
the physics of the plane thickness (P1) and at detector level with pile-up 200 for a search and the measurement of $\mQ$
and $\Lam$ (P2). The study uses public information only and is not a CMS result.

\paragraph{P1 verdict (truth level).} The B-field thickness of the quirk plane is generated at the near-miss crossings of the
oscillating pair: the solenoid field, an electric field in the pair frame, makes the ends miss by $\bodd$ at every second
crossing, and the string's transverse pull integrates to $\ln(\lmax/\bodd)$, which rotates the oscillation axis. This gives
an analytic form of the per-period rotation that Li et al.\ determined numerically, and the law
$\DsB\propto\epsB[\ln(1/\epsB)+C]$ with local slope $c=1/[\ln(1/\epsB)+C]\simeq0.104$ at $1\keV$ and $500\GeV$, i.e.\
$\DsB\approx\Lam^{-1.8}B^{0.9}$, nearly independent of $\mQ$. The law is confirmed by blind numerics at $1\keV$ and holds per
event on $2.2\times10^5$ simulated events within $\lmax/\RIT\lesssim0.6$ and $\kB\lesssim4$; the population median of $c$
exceeds the prediction by $2$--$3\%$ (a failed pre-set gate, read as an estimator effect), a larger $\kB$-dependent excess at
$3$--$5\keV$ is attributed to the CM magnetic field as an estimate, and one prediction (the straight-string ratio at $\gst=2$)
fails. Multiple scattering sets a floor of 10--28$\um$ and dominates above $\Lcross\approx1.2$--$1.9\keV$, and the two add
below quadrature in the crossover. For a search, the field thickness matters below $\approx1.5\keV$; above it, the finder
tolerance is set by material.

\paragraph{P2 verdict (detector level).} With a dedicated barrel plane finder, a pile-up overlay normalised to the TDR
occupancies and a working-point rule frozen before any signal efficiency, \textbf{no working point exists}: the lowest
background is $\Nbkgeightyfour=1.23\times10^6$ events (R09), 7.05 decades above the budget, and the deep-rung background
consists of zero-count bounds. The pre-registered fallback gives expected limits at best 790 times the predicted signal
($E$ 200, $1.5\keV$; $\approx910$--950 with the full-chain efficiency systematic; ratios corrected for the
$p_T(Q\bar Q)>200\GeV$ tail fraction of the cross section, which an earlier version omitted), so \textbf{nothing is excluded}, and the
study adds nothing to the existing bounds (all below $\approx180\GeV$ in this $\Lam$ range). At pile-up 200 neither $\mQ$ nor
$\Lam$ is measured: the time-of-flight estimator is a pair-mass proxy with a pile-up-induced low cluster, and the
$\Lam^{-2}$ separation handle is not measurable. The negative result is dominated by simulation statistics and by the
pile-up extrapolation uncertainty, not by an observed background.

\paragraph{What would be needed.} None of the following can be applied to the frozen analysis; each needs a new
pre-registration. (i) A \emph{data-driven pile-up fake estimate} in control regions (MET-road or $\Delta\phi$ sidebands, or
out-of-time crossings) to replace the thinning extrapolation and its trend and closure uncertainties, the decisive item.
(ii) Deeper pile-up simulation ($\rho=3$, or a closure test with a real lever arm at R07--R09), worth about one decade.
(iii) Single-interaction conditionals from the high-statistics no-pile-up gun samples with their own closure test (about two
decades on those bounds) and time-tag credit for prompt SM and single-interaction terms in SR-T (up to four decades there).
(iv) A merging-tolerant counting rung or a line-topology category at $\Lam\gtrsim3\keV$, a lower seed-opening cut and an
endcap finder, to raise the $\approx0.3$ barrel ceiling of the efficiency. (v) A pile-up-robust separation estimator and
in-time-hit rejection for the time of flight, to recover the $\Lam$ and $\mQ$ handles. Even with (ii) and (iii) the estimate
leaves $\approx0$--2 (SR-T) and 4--5 (SR-G) decades (\EST{}, from the previous analysis run; order of magnitude
unchanged), so (i) is essential. Validating the background inputs with
CMS Open Data (Section~\ref{sec:opendata}) and an analytic multiple-scattering theory of the plane thickness are natural
follow-ups.
